\chapter{Principal and Frame Bundles}
\label{ch:vii08-principal-and-frame-bundles}

Vector bundles encode linear fibers. Principal bundles isolate the symmetry behind their changes of frame. A principal \(G\)-bundle has fibers that look like the group \(G\) itself, but with no preferred identity element in an individual fiber.
\[
\boxed{P\times G\to P,\qquad (u,g)\mapsto u\cdot g,\qquad
P|_U\cong U\times G.}
\]

\section{Smooth group actions}
\begin{definition}[Right action]\label{def:vii08-right-action}
A smooth right action of a Lie group \(G\) on a manifold \(P\) is a smooth map \(P\times G\to P\), \((u,g)\mapsto u\cdot g\), such that \(u\cdot e=u\) and \((u\cdot g)\cdot h=u\cdot(gh)\).
\end{definition}

\begin{definition}[Free and transitive]\label{def:vii08-free-transitive}
An action is free if \(u\cdot g=u\) implies \(g=e\). It is transitive on a subset if any two points of that subset differ by the action of some group element.
\end{definition}

\section{Principal bundles}
\begin{definition}[Principal \(G\)-bundle]\label{def:vii08-principal-bundle}
A principal \(G\)-bundle over \(M\) is a smooth manifold \(P\), a smooth surjection \(\pi:P\to M\), and a smooth right \(G\)-action preserving fibers, such that every \(p\in M\) has a neighborhood \(U\) and a \(G\)-equivariant diffeomorphism
\[
\Phi:\pi^{-1}(U)\to U\times G,
\qquad
\Phi(u\cdot g)=(p,hg)\quad\text{if }\Phi(u)=(p,h).
\]
\end{definition}

\begin{proposition}[Fiber torsor property]\label{prop:vii08-torsor}
The \(G\)-action on every fiber \(P_p\) is free and transitive. Hence for \(u,v\in P_p\) there is a unique \(g\in G\) with \(v=u\cdot g\).
\end{proposition}

\begin{example}[Trivial principal bundle]\label{ex:vii08-trivial}
The projection \(M\times G\to M\) with \((p,h)\cdot g=(p,hg)\) is the trivial principal \(G\)-bundle.
\end{example}

\section{Local sections and transition functions}
\begin{proposition}[Local section versus trivialization]\label{prop:vii08-local-section}
A local section \(s:U\to P\) determines a trivialization by
\[
U\times G\to P|_U,\qquad (p,g)\mapsto s(p)\cdot g,
\]
and every principal-bundle trivialization determines a local section by taking the group identity.
\end{proposition}

\begin{theorem}[Principal transition functions]\label{thm:vii08-transition}
For local sections \(s_\alpha,s_\beta\), there is a unique smooth map
\[
g_{\alpha\beta}:U_\alpha\cap U_\beta\to G
\]
such that \(s_\beta=s_\alpha\cdot g_{\alpha\beta}\). These satisfy
\[
g_{\alpha\alpha}=e,\qquad g_{\beta\alpha}=g_{\alpha\beta}^{-1},\qquad
g_{\alpha\gamma}=g_{\alpha\beta}g_{\beta\gamma}.
\]
\end{theorem}

\begin{warning}[Convention matters]\label{warning:vii08-convention}
Because we use a right group action, the order in the cocycle identity is tied to the convention \(s_\beta=s_\alpha g_{\alpha\beta}\). Switching to left actions changes formulas. One should fix a convention before manipulating transition functions.
\end{warning}

\section{The frame bundle of a vector bundle}
\begin{definition}[Frame]\label{def:vii08-frame}
If \(E\to M\) has rank \(r\), a frame at \(p\) is an ordered basis \((e_1,\ldots,e_r)\) of \(E_p\), equivalently a linear isomorphism
\[
u:\mathbb R^r\to E_p.
\]
\end{definition}

\begin{definition}[Frame bundle]\label{def:vii08-frame-bundle}
The frame bundle \(\operatorname{Fr}(E)\) is the set of all frames in all fibers. It carries a right \(GL(r,\mathbb R)\)-action
\[
u\cdot A=u\circ A,
\]
and projection to the base point of the frame.
\end{definition}

\begin{theorem}[Frame bundle is principal]\label{thm:vii08-frame-principal}
For a rank-\(r\) vector bundle \(E\to M\), \(\operatorname{Fr}(E)\to M\) is a principal \(GL(r,\mathbb R)\)-bundle.
\end{theorem}
\begin{proof}
A bundle frame \(e_1,\ldots,e_r\) over \(U\) identifies every frame over \(p\in U\) uniquely with a matrix \(A\in GL(r)\). This gives an equivariant diffeomorphism \(\operatorname{Fr}(E)|_U\cong U\times GL(r)\).
\end{proof}

\begin{corollary}[Linear frame bundle of a manifold]\label{cor:vii08-lin-frame}
For an \(n\)-manifold \(M\), \(\operatorname{Fr}(TM)\) is a principal \(GL(n,\mathbb R)\)-bundle, called the linear frame bundle of \(M\).
\end{corollary}

\section{Principal bundle maps}
\begin{definition}[Principal-bundle morphism]\label{def:vii08-principal-map}
A morphism of principal \(G\)-bundles over \(f:M\to N\) is a smooth map \(F:P\to Q\) satisfying \(\pi_QF=f\pi_P\) and \(F(u\cdot g)=F(u)\cdot g\).
\end{definition}

\begin{proposition}[Equivariant maps are determined fiberwise]\label{prop:vii08-equivariant-determined}
If \(F:P\to Q\) is equivariant and its value is known at one point of a fiber, its value on the entire fiber is forced by equivariance.
\end{proposition}

\section{Associated bundles}
Let \(\rho:G\to GL(V)\) be a finite-dimensional representation.
\begin{definition}[Associated vector bundle]\label{def:vii08-associated}
For a principal \(G\)-bundle \(P\to M\), define
\[
P\times_\rho V=(P\times V)/\sim,
\qquad
(u\cdot g,v)\sim(u,\rho(g)v).
\]
The projection \([u,v]\mapsto\pi(u)\) makes this a vector bundle over \(M\).
\end{definition}

\begin{theorem}[Recovering a vector bundle from its frames]\label{thm:vii08-recover-vector}
If \(E\to M\) has rank \(r\) and \(\rho\) is the standard representation of \(GL(r)\) on \(\mathbb R^r\), then
\[
\operatorname{Fr}(E)\times_{GL(r)}\mathbb R^r\cong E,
\qquad [u,v]\mapsto u(v).
\]
\end{theorem}
\begin{proof}
The map is well defined because \((uA,A^{-1}v)\) and \((u,v)\) represent the same geometric vector under the consistent associated-bundle convention. Fiberwise it is a linear bijection, and local frame trivializations show smoothness.
\end{proof}

\begin{remark}[Convention for associated bundles]\label{rem:vii08-associated-convention}
Equivalent conventions place \(g\) or \(g^{-1}\) in the relation. What matters is consistency with the chosen right action and representation convention.
\end{remark}

\section{Gauge changes without connections}
\begin{definition}[Gauge transformation in a trivialization]\label{def:vii08-gauge-change}
Changing a local section from \(s\) to \(s'=s\cdot h\), where \(h:U\to G\) is smooth, changes the local principal-bundle coordinates. This is a local gauge change.
\end{definition}

\begin{proposition}[Transition functions under gauge change]\label{prop:vii08-gauge-transition}
If \(s'_\alpha=s_\alpha h_\alpha\), then the transition functions change by
\[
g'_{\alpha\beta}=h_\alpha^{-1}g_{\alpha\beta}h_\beta.
\]
\end{proposition}
\begin{proof}
Substitute the changed sections into \(s'_\beta=s'_\alpha g'_{\alpha\beta}\) and use freeness of the right action.
\end{proof}

\section{Examples and geometric meaning}
\begin{example}[Frames on a trivial bundle]\label{ex:vii08-trivial-frame}
For \(E=M\times\mathbb R^r\), the frame bundle is canonically \(M\times GL(r,\mathbb R)\).
\end{example}

\begin{example}[Frame transitions from bundle transitions]\label{ex:vii08-frame-transition}
If a vector bundle has transition matrices \(g_{\beta\alpha}\), the induced frame-bundle transition functions are the same \(GL(r)\)-valued maps, interpreted as changes of ordered basis.
\end{example}

\begin{warning}[Connections come later]\label{warning:vii08-no-connection}
A principal bundle supplies fibers, symmetry, local sections, and transition functions. It does not yet choose horizontal directions or parallel transport. Those additional structures belong to VII/21 and later chapters.
\end{warning}

\section{Exercises}

\begin{exercise}\label{exr:vii08-01}
Verify that right multiplication defines a smooth right action of \(G\) on itself.
\end{exercise}
\begin{hint}
Check identity and associativity.
\end{hint}
\begin{solution}
The group multiplication map is smooth, \(h e=h\), and \((hg)k=h(gk)\).
\end{solution}

\begin{exercise}\label{exr:vii08-02}
Show that the action on a principal-bundle fiber is free.
\end{exercise}
\begin{hint}
Use a local trivialization.
\end{hint}
\begin{solution}
In \(U\times G\), \((p,h)g=(p,h)\) implies \(hg=h\), hence \(g=e\).
\end{solution}

\begin{exercise}\label{exr:vii08-03}
Show that the action on a principal-bundle fiber is transitive.
\end{exercise}
\begin{hint}
Again use \(U\times G\).
\end{hint}
\begin{solution}
Given \((p,h_1)\) and \((p,h_2)\), choose \(g=h_1^{-1}h_2\).
\end{solution}

\begin{exercise}\label{exr:vii08-04}
Prove that a local section produces a principal-bundle trivialization.
\end{exercise}
\begin{hint}
Use \((p,g)\mapsto s(p)g\).
\end{hint}
\begin{solution}
Freeness and transitivity give a bijection on each fiber; local smoothness and the smooth division map in a principal chart give a diffeomorphism.
\end{solution}

\begin{exercise}\label{exr:vii08-05}
Recover a local section from a trivialization.
\end{exercise}
\begin{hint}
Insert the identity element.
\end{hint}
\begin{solution}
Define \(s(p)=\Phi^{-1}(p,e)\); this is smooth and projects to \(p\).
\end{solution}

\begin{exercise}\label{exr:vii08-06}
Derive the transition function relating two local sections.
\end{exercise}
\begin{hint}
Use the torsor property.
\end{hint}
\begin{solution}
For each \(p\), there is a unique \(g_{\alpha\beta}(p)\) with \(s_\beta(p)=s_\alpha(p)g_{\alpha\beta}(p)\); local principal coordinates show smoothness.
\end{solution}

\begin{exercise}\label{exr:vii08-07}
Prove the principal cocycle identity.
\end{exercise}
\begin{hint}
Compare \(s_\gamma\) in two ways.
\end{hint}
\begin{solution}
From \(s_\gamma=s_\beta g_{\beta\gamma}=s_\alpha g_{\alpha\beta}g_{\beta\gamma}\) and freeness, \(g_{\alpha\gamma}=g_{\alpha\beta}g_{\beta\gamma}\).
\end{solution}

\begin{exercise}\label{exr:vii08-08}
Show that \(M\times G\to M\) has the global section \(p\mapsto(p,e)\).
\end{exercise}
\begin{hint}
Use the identity of \(G\).
\end{hint}
\begin{solution}
The map is smooth and projects to \(p\).
\end{solution}

\begin{exercise}\label{exr:vii08-09}
Show that a principal bundle with a global section is trivial.
\end{exercise}
\begin{hint}
Use the section-to-trivialization construction globally.
\end{hint}
\begin{solution}
The map \((p,g)\mapsto s(p)g\) is a global equivariant diffeomorphism \(M\times G\to P\).
\end{solution}

\begin{exercise}\label{exr:vii08-10}
Show that a frame is equivalently a linear isomorphism \(\mathbb R^r\to E_p\).
\end{exercise}
\begin{hint}
Send the standard basis to the ordered basis.
\end{hint}
\begin{solution}
Every ordered basis determines a unique linear isomorphism, and conversely the images of the standard basis form a frame.
\end{solution}

\begin{exercise}\label{exr:vii08-11}
Verify the right \(GL(r)\)-action on frames is free.
\end{exercise}
\begin{hint}
Use injectivity of a frame map.
\end{hint}
\begin{solution}
If \(uA=u\), then composing with \(u^{-1}\) gives \(A=I\).
\end{solution}

\begin{exercise}\label{exr:vii08-12}
Verify the frame action is transitive on a fiber.
\end{exercise}
\begin{hint}
Compare two isomorphisms \(u,v:\mathbb R^r\to E_p\).
\end{hint}
\begin{solution}
The unique matrix is \(A=u^{-1}v\), so \(v=uA\).
\end{solution}

\begin{exercise}\label{exr:vii08-13}
Compute \(\dim \operatorname{Fr}(E)\) for a rank-\(r\) bundle over an \(n\)-manifold.
\end{exercise}
\begin{hint}
Use local products with \(GL(r)\).
\end{hint}
\begin{solution}
Since \(\dim GL(r)=r^2\), the total dimension is \(n+r^2\).
\end{solution}

\begin{exercise}\label{exr:vii08-14}
Show \(\operatorname{Fr}(M\times\mathbb R^r)\cong M\times GL(r)\).
\end{exercise}
\begin{hint}
Compare every frame to the standard frame.
\end{hint}
\begin{solution}
A frame is uniquely a standard frame followed by some \(A\in GL(r)\), smoothly in \(p\).
\end{solution}

\begin{exercise}\label{exr:vii08-15}
Explain why \(\operatorname{Fr}(TM)\) has structure group \(GL(n)\).
\end{exercise}
\begin{hint}
A tangent frame is a basis of an \(n\)-dimensional vector space.
\end{hint}
\begin{solution}
The change from one ordered tangent basis to another is a unique element of \(GL(n,\mathbb R)\).
\end{solution}

\begin{exercise}\label{exr:vii08-16}
Check that a principal-bundle morphism preserves the group coordinate in local trivializations.
\end{exercise}
\begin{hint}
Use equivariance.
\end{hint}
\begin{solution}
Once the image of a chosen local section is known, equivariance forces \(F(s(p)g)=F(s(p))g\).
\end{solution}

\begin{exercise}\label{exr:vii08-17}
Show the associated-bundle equivalence relation is compatible with the group action convention.
\end{exercise}
\begin{hint}
Check composition by two group elements.
\end{hint}
\begin{solution}
Applying \(g\) then \(h\) produces the same class as applying \(gh\), precisely because \(\rho\) is a representation.
\end{solution}

\begin{exercise}\label{exr:vii08-18}
Compute the rank of \(P\times_\rho V\).
\end{exercise}
\begin{hint}
Use a principal local trivialization.
\end{hint}
\begin{solution}
Over \(U\), choosing a local section identifies the associated bundle with \(U\times V\), so its rank is \(\dim V\).
\end{solution}

\begin{exercise}\label{exr:vii08-19}
Show that the standard associated bundle of \(\operatorname{Fr}(E)\) has fiber \(E_p\).
\end{exercise}
\begin{hint}
Evaluate a frame on a vector.
\end{hint}
\begin{solution}
The map \([u,v]\mapsto u(v)\) is fiberwise a well-defined linear bijection.
\end{solution}

\begin{exercise}\label{exr:vii08-20}
Derive the gauge-transformation law for principal transition functions.
\end{exercise}
\begin{hint}
Insert \(s'_\alpha=s_\alpha h_\alpha\).
\end{hint}
\begin{solution}
Substitution into the chosen right-action convention gives \(g'_{\alpha\beta}=h_{\alpha}^{-1}g_{\alpha\beta}h_{\beta}\).
\end{solution}

\begin{exercise}\label{exr:vii08-21}
For an abelian group \(G\), simplify the gauge-transformation law.
\end{exercise}
\begin{hint}
Commute the factors.
\end{hint}
\begin{solution}
It becomes \(g'_{\alpha\beta}=g_{\alpha\beta}h_\alpha^{-1}h_\beta\).
\end{solution}

\begin{exercise}\label{exr:vii08-22}
Show that a global frame of \(E\) is the same as a global section of \(\operatorname{Fr}(E)\).
\end{exercise}
\begin{hint}
A point of the frame bundle is itself a frame.
\end{hint}
\begin{solution}
A section selects smoothly one frame at every base point, exactly a global frame of \(E\).
\end{solution}

\begin{exercise}\label{exr:vii08-23}
Use the previous exercise to restate the triviality criterion for vector bundles in principal-bundle language.
\end{exercise}
\begin{hint}
Apply the global-section criterion.
\end{hint}
\begin{solution}
\(E\) is trivial iff \(\operatorname{Fr}(E)\) has a global section, equivalently iff the frame principal bundle is trivial.
\end{solution}

\begin{exercise}\label{exr:vii08-24}
State what geometric datum has deliberately not been added in this chapter.
\end{exercise}
\begin{hint}
Read the final warning.
\end{hint}
\begin{solution}
No horizontal distribution, connection form, covariant derivative, or parallel transport has been chosen.
\end{solution}


\section{Solved problem dossiers}

\begin{problem}[Legacy principal-bundle problem: torsors]\label{prob:vii08-01}
Prove intrinsically that each principal-bundle fiber is a \(G\)-torsor.
\end{problem}
\begin{solution}
A local principal trivialization identifies the fiber with \(G\) carrying right multiplication. Right multiplication is free and transitive, and these properties are invariant under equivariant diffeomorphism.
\end{solution}

\begin{problem}[Legacy local-section problem]\label{prob:vii08-02}
Prove the equivalence between local sections and local principal trivializations.
\end{problem}
\begin{solution}
A section \(s\) gives \((p,g)\mapsto s(p)g\). Conversely a trivialization gives \(s(p)=\Phi^{-1}(p,e)\). These constructions are inverse to one another.
\end{solution}

\begin{problem}[Legacy transition-function problem]\label{prob:vii08-03}
Derive transition functions from local sections and prove the cocycle condition.
\end{problem}
\begin{solution}
Use the unique group element relating two section values in each fiber. Smoothness follows in a principal chart; comparing three sections gives the cocycle law.
\end{solution}

\begin{problem}[Legacy frame-bundle problem]\label{prob:vii08-04}
Construct the principal \(GL(r)\)-bundle structure on \(\operatorname{Fr}(E)\).
\end{problem}
\begin{solution}
A local frame identifies any frame uniquely with a matrix in \(GL(r)\), giving \(\operatorname{Fr}(E)|_U\cong U\times GL(r)\). Composition on the right is exactly the principal action.
\end{solution}

\begin{problem}[Tangent-frame dossier]\label{prob:vii08-05}
Describe the frame bundle of an \(n\)-manifold directly in coordinates.
\end{problem}
\begin{solution}
Over a coordinate chart, the coordinate tangent frame is a local section of \(\operatorname{Fr}(TM)\). Every other tangent frame is obtained by a unique \(GL(n)\) matrix; chart changes induce Jacobian-valued transition functions.
\end{solution}

\begin{problem}[Global-section dossier]\label{prob:vii08-06}
Prove a principal \(G\)-bundle is trivial iff it has a global section.
\end{problem}
\begin{solution}
A trivial bundle has the identity section. A global section gives the global equivariant trivialization \((p,g)\mapsto s(p)g\).
\end{solution}

\begin{problem}[Associated-bundle dossier]\label{prob:vii08-07}
Construct the local trivializations of \(P\times_\rho V\).
\end{problem}
\begin{solution}
Choose a local section \(s\). Every \(u\in P_p\) is uniquely \(s(p)g\); rewrite \([u,v]\) using the equivalence relation as a class represented by \([s(p),w]\). This identifies the fiber with \(V\), smoothly over \(U\).
\end{solution}

\begin{problem}[Recovering \(E\) from frames]\label{prob:vii08-08}
Prove \(\operatorname{Fr}(E)\times_{GL(r)}\mathbb R^r\cong E\).
\end{problem}
\begin{solution}
Evaluation \([u,v]\mapsto u(v)\) is compatible with the quotient relation, linear and bijective in every fiber, and in a local frame it is the identity trivialization.
\end{solution}

\begin{problem}[Gauge-change dossier]\label{prob:vii08-09}
Show how changing local sections conjugates transition functions.
\end{problem}
\begin{solution}
Write \(s'_\alpha=s_\alpha h_\alpha\), impose \(s'_\beta=s'_\alpha g'_{\alpha\beta}\), and cancel \(s_\alpha\) using freeness to obtain \(g'_{\alpha\beta}=h_\alpha^{-1}g_{\alpha\beta}h_\beta\).
\end{solution}

\begin{problem}[Equivariant-map dossier]\label{prob:vii08-10}
Show an equivariant map between principal bundles is determined by its values on a local section.
\end{problem}
\begin{solution}
Every point over \(U\) is uniquely \(s(p)g\). Equivariance forces \(F(s(p)g)=F(s(p))g\), so no additional choice remains.
\end{solution}

\begin{problem}[Vector/principal dictionary]\label{prob:vii08-11}
Match vector-bundle frames, principal local sections, and vector-bundle local trivializations.
\end{problem}
\begin{solution}
A vector local frame is a section of the frame bundle. Such a principal section trivializes the frame bundle and simultaneously identifies every vector in \(E_p\) with its coordinate column in \(\mathbb R^r\).
\end{solution}

\begin{problem}[Why no preferred identity in a fiber]\label{prob:vii08-12}
Explain why a principal fiber resembles \(G\) but is not canonically the group \(G\).
\end{problem}
\begin{solution}
Choosing one point \(u\in P_p\) identifies \(g\mapsto ug\) with \(G\), but another choice translates this identification. Without a chosen section there is no canonical identity element in the fiber.
\end{solution}


\section{Challenge problems}

\begin{problem}[Reconstruct a principal bundle from a cocycle]\label{prob:vii08-ch01}
Given smooth \(G\)-valued transition functions satisfying the cocycle law, construct a principal \(G\)-bundle.
\end{problem}
\begin{solution}
Glue \(U_\alpha\times G\) with the equivalence determined by the transition functions. The cocycle law makes the quotient consistent; right multiplication descends because the gluing maps are equivariant; the product pieces become local principal trivializations.
\end{solution}

\begin{problem}[Associated bundle functoriality]\label{prob:vii08-ch02}
Let \(T:V\to W\) intertwine two representations \(\rho\) and \(\sigma\). Show it induces a vector-bundle map \(P\times_\rho V\to P\times_\sigma W\).
\end{problem}
\begin{solution}
Define \([u,v]\mapsto[u,Tv]\). If representatives are changed by \(g\), equivariance \(T\rho(g)=\sigma(g)T\) ensures the resulting class is unchanged. The map is locally \(\operatorname{id}_U\times T\), hence smooth.
\end{solution}

\begin{problem}[Frame bundle and bundle isomorphisms]\label{prob:vii08-ch03}
Show a vector-bundle isomorphism \(F:E\to E'\) over \(M\) induces a principal-bundle isomorphism \(\operatorname{Fr}(E)\to\operatorname{Fr}(E')\).
\end{problem}
\begin{solution}
Send a frame \(u:\mathbb R^r\to E_p\) to \(F_p\circ u\). This is smooth in local frames and commutes with the right \(GL(r)\)-action: \(F_p\circ(uA)=(F_p\circ u)A\).
\end{solution}

\begin{problem}[Triviality equivalence]\label{prob:vii08-ch04}
Prove carefully that \(E\) is trivial iff its frame bundle is a trivial principal \(GL(r)\)-bundle.
\end{problem}
\begin{solution}
A global vector frame is exactly a global section of \(\operatorname{Fr}(E)\). VII/07 says a global frame trivializes \(E\); the principal global-section criterion says the same section trivializes \(\operatorname{Fr}(E)\).
\end{solution}

\begin{problem}[Why connections are additional structure]\label{prob:vii08-ch05}
Using only local sections and transition functions, explain why no canonical horizontal complement to the vertical tangent spaces has yet appeared.
\end{problem}
\begin{solution}
The principal-bundle axioms determine the vertical directions generated by the group action, but local gauge changes can mix any proposed complementary directions unless extra transformation-compatible data are specified. A connection is precisely such an additional choice, postponed to VII/21.
\end{solution}

\section*{Bridge to VII/09}
Principal bundles organize changes of frame. Differential forms, by contrast, are tensorial fields that can be pulled back along arbitrary smooth maps without choosing a frame. VII/09 develops alternating covariant tensors, wedge products, pullbacks, and the exterior derivative.
